\section{АРХИТЕКТУРА СИСТЕМЫ УПРАВЛЕНИЯ СЧЕТАМИ}

\subsection{Компонентная структура и распределение ролей}
Архитектура информационной системы \texttt{invoices-s13} базируется на паттерне <<API Gateway>>, который является стандартом для современных микросервисных приложений. Система разделена на два функциональных контура:

\begin{enumerate}
    \item \textbf{API Gateway (Внешний контур)}: Выступает в роли единой точки входа (Single Point of Entry). Его основные задачи включают маршрутизацию запросов, агрегацию данных из различных сервисов (в перспективе), выполнение функций обратного прокси и трансляцию протоколов (Protocol Translation) из внешнего REST/JSON во внутренний gRPC/Protobuf. Использование шлюза позволяет скрыть детали реализации бэкенда от клиента.
    \item \textbf{Invoices Service (Внутренний контур)}: Является <<сервисом-источником истины>> для финансовых документов. Он инкапсулирует всю бизнес-логику: валидацию сумм (\texttt{amount}), генерацию идентификаторов и управление состоянием счетов. Сервис спроектирован так, чтобы быть максимально изолированным и независимым от способа отображения данных клиенту.
\end{enumerate}

\subsection{Стек протоколов и обоснование выбора}
Для обеспечения баланса между доступностью и производительностью в системе реализован гибридный подход к коммуникациям:

\begin{itemize}
    \item \textbf{REST (HTTP/1.1 + JSON)}: Используется для взаимодействия <<клиент-шлюз>>. Выбор обусловлен широкой поддержкой протокола всеми современными языками программирования и браузерами. JSON обеспечивает легкость отладки и читаемость сообщений.
    \item \textbf{gRPC (HTTP/2 + Protobuf)}: Используется для взаимодействия <<шлюз-сервис>>. В отличие от REST, gRPC использует бинарную сериализацию Protobuf, что существенно (в 3-5 раз) уменьшает размер передаваемого трафика. Использование HTTP/2 позволяет мультиплексировать несколько запросов в одном TCP-соединении, что критично для высоконагруженных распределенных систем.
\end{itemize}

\subsection{Модель управления данными и жизненный цикл}
В текущей реализации используется оперативная память (in-memory) для хранения объектов \texttt{invoice}. Это решение позволяет минимизировать время отклика системы на этапе тестирования. Данные организованы в виде потокобезопасного словаря, где ключом выступает строковый ID.

Архитектура системы поддерживает горизонтальное масштабирование: при увеличении нагрузки можно запустить несколько экземпляров \texttt{invoices-svc-s13}. Однако, для корректной работы в таком режиме потребуется внедрение распределенной базы данных (например, PostgreSQL или Redis), что предусмотрено интерфейсом репозитория в коде.

\subsection{Оркестрация и Service Discovery}
Взаимодействие компонентов в динамической среде Docker обеспечивается встроенным механизмом Service Discovery. Шлюз не использует статические IP-адреса, а обращается к бэкенду по его логическому имени \texttt{invoices-svc-s13}. Внутренняя DNS-служба Docker разрешает это имя в актуальный IP-адрес контейнера, что позволяет прозрачно обновлять и перезапускать сервисы.
